\documentclass[11pt]{article}
\usepackage[margin=1in]{geometry}
\usepackage{times}
\usepackage{hyperref}
\hypersetup{colorlinks=true, linkcolor=black, urlcolor=blue, citecolor=black}
\title{Spin Without a Source: Torsion Field Claims as a Gravitational Mechanism}
\author{Flyxion \\ \small Independent Researcher}
\date{}
\begin{document}
\maketitle

\begin{abstract}
Torsion field theories, most prominently associated with the work of Anatoly Akimov and Gennady Shipov, propose a subtle, information-carrying field generated by rotating or spinning matter, distinct from gravitation and electromagnetism, capable of instantaneous transmission and, in some formulations, of locally altering weight or inertia. This paper distinguishes the legitimate mathematical object of torsion in differential geometry, a genuine and well-understood generalization of Riemannian curvature used in some extensions of general relativity, from the informal, empirically unconstrained torsion field of the Akimov-Shipov tradition, and argues that the latter borrows the former's name and superficial vocabulary while dispensing with everything that made it mathematically or observationally accountable. The paper recounts how the claims were funded and reviewed in Russia, applies the standard coupling of torsion to spin to show how small the legitimate effect is expected to be, and examines the defense that the field carries information without energy. It makes no claim about the sincerity of any proponent.
\end{abstract}

\section{Introduction}

Torsion, in the technical sense used in Einstein-Cartan theory and related extensions of general relativity, is the antisymmetric part of a connection's coefficients, a well-defined geometric object that couples to the intrinsic spin density of matter rather than to its mass-energy alone. The idea goes back to Cartan in the 1920s and was developed into a gravitational theory with spin by Sciama and by Kibble \cite{K1961}, and the standard review of the resulting framework is that of Hehl, von der Heyde, Kerlick and Nester \cite{HHKN1976}. The effects of torsion in this framework are exceedingly small under ordinary conditions, because the torsion is sourced algebraically by spin density, with the gravitational constant as the coupling, and a modest macroscopic spin density produces a torsion far too small for any laboratory to see. They have never been detected, consistent with theoretical expectation, at any laboratory or astrophysical scale probed so far, and the question of whether and how they could be detected is reviewed in the literature on the physical aspects of space-time torsion \cite{S2002}. This is a legitimate and actively studied area of theoretical physics, if a comparatively minor one.

The torsion field of Akimov and Shipov's popular and commercial work, sometimes called spin field or axion field in loosely translated Russian-language literature from the 1990s, shares essentially nothing with this except the word torsion. It is described as generated by any rotating object, propagating without attenuation or delay across arbitrary distances, carrying information rather than energy, and capable of influencing biological systems, consciousness, and the weight of nearby objects. None of these properties follows from, or is consistent with, the differential-geometric object the name is borrowed from. Shipov's Theory of the Physical Vacuum presents the field as an independent physical entity analogous to the electromagnetic field, generated by spin rather than by charge \cite{AS}.

\section{How the Claims Were Funded and Reviewed}

The history matters because it shows that the claims were not ignored by the scientific community but examined and found wanting. Akimov directed a state-sponsored center for non-traditional technologies in the late Soviet period, and accounts attributed to the Russian Academy of Sciences state that large sums were spent on torsion research without scientific examination until an inquiry by the Academy in 1991. Funding from the Russian Ministry of Science and later the Ministry of Defense is reported for the 1990s. The Commission against Pseudoscience and Falsification of Scientific Research was established within the Russian Academy of Sciences in 1998 and 1999, and its chairman named the torsion field program among the examples of pseudoscientific projects that had obtained state money \cite{Kr2002}. A device built on torsion technologies was installed in a Russian satellite launched in 2008, and the scientists who reviewed its performance concluded that it had no effect on the satellite's orbit \cite{MT}.

I cite these accounts as a description of how the claims fared under institutional review, not as an argument that the claims are false because a commission said so. The substantive argument follows from the physics.

\section{Why "Carries Information Without Carrying Energy" Is Not a Loophole}

A recurring defense of the informal torsion field is that its undetectability by conventional instruments is expected, since it is claimed to carry information rather than energy, and instruments built to detect energy transfer would naturally miss a signal of this kind. This framing treats the absence of energy transfer as a feature that explains away the absence of detection, but information transfer without any accompanying physical interaction is not a subtle experimental challenge, it is a description of something that cannot interact with any detector at all, by definition, since detection is itself a physical interaction. A field that in principle cannot be measured by any physical process is not describing a hidden layer of physics; it is describing the absence of a testable claim.

The same difficulty appears in the claim of instantaneous transmission. A field that propagates with infinite speed is incompatible with relativistic causality in the standard framework, which is not a modest correction to be accommodated but a structural feature of every field theory consistent with the observed behavior of light and matter. The proponents' answer has been that the torsion field lies outside the relativistic framework, a position that removes the claim from the body of physics that its vocabulary borrows from, while still using the language of that body to give the claim an appearance of continuity with it.

\section{The Entropic Gravity Framing}

On an account where gravitational curvature tracks the redistribution of an ordered configuration across coupled scalar, vector, and entropy fields, weight and inertia are not independently adjustable quantities that a separate, parallel field could tune from outside; they are expressions of how a given mass-energy distribution participates in that redistribution. A device claiming to alter local weight through spin-generated torsion fields would need a mechanism by which ordinary mechanical rotation, itself an entirely mundane redistribution of angular momentum with no unusual entropy signature, somehow taps into and modifies the much larger-scale bookkeeping that produces gravitational curvature. No such channel exists in this framework, any more than in standard general relativity, and the informal torsion field literature has never proposed one with enough specificity to evaluate, relying instead on qualitative language about subtle energies and information fields that resists translation into any testable prediction.

\section{Objections}

Proponents sometimes cite anecdotal or small-sample biological and materials-science experiments reporting effects from torsion generators, including changes in crystal structure, seed germination rates, or reaction times in human subjects. These studies, where they exist in published form, are characterized by small sample sizes, absence of blinding, and no independent replication by laboratories outside the group that produced the original device, a pattern indistinguishable from publication bias and experimenter effect rather than a genuine physical signal.

A second objection holds that torsion is a real quantity in theoretical physics, and that to dismiss torsion fields is to dismiss a respectable subject. This is a confusion between a name and a thing, and it is the confusion the paper is about. The respectable subject has definite equations, a definite coupling and a definite prediction that its effects are negligible at laboratory scales. The informal field has none of these. Pointing to the former does not support the latter.

\section{Conclusion}

The differential-geometric object called torsion is real, well defined, and predicted to produce effects far too small to explain any of the claims made under its borrowed name. The Akimov-Shipov torsion field, evaluated on its own terms, describes an interaction with no available mechanism to alter gravitational curvature and no falsifiable prediction that has survived independent scrutiny, and the informational framing offered in its defense removes the claim from the category of physics rather than protecting it within it. This paper has not tested a torsion generator and does not deny that some effects reported in the literature of such devices may have conventional explanations of their own.

\begin{thebibliography}{99}

\bibitem{K1961} T. W. B. Kibble, Lorentz Invariance and the Gravitational Field, \emph{Journal of Mathematical Physics} 2 (1961) 212--221.

\bibitem{HHKN1976} F. W. Hehl, P. von der Heyde, G. D. Kerlick and J. M. Nester, General Relativity with Spin and Torsion: Foundations and Prospects, \emph{Reviews of Modern Physics} 48 (1976) 393--416.

\bibitem{S2002} I. L. Shapiro, Physical Aspects of the Space-Time Torsion, \emph{Physics Reports} 357 (2002) 113--213, arXiv:hep-th/0103093.

\bibitem{AS} A. E. Akimov, article on torsion fields, \emph{New Energy Technologies}, Issue 1 (16), January--March 2004, and G. I. Shipov, \emph{A Theory of Physical Vacuum}, Moscow.

\bibitem{Kr2002} E. Kruglyakov, Why Is Pseudoscience Dangerous?, \emph{Skeptical Inquirer}, Center for Inquiry, 2002.

\bibitem{MT} Interview on pseudoscience in Russia, \emph{The Moscow Times}, with a note on the 2008 Yubileiny satellite torsion engine.

\end{thebibliography}

\end{document}
